\section{Verification}\label{sec:verification}

Every obligation of a \VG{} node becomes one quantifier-free SMT query: the conjunction of
some assumptions with the negation of a goal. An unsatisfiable query proves the obligation; a
satisfiable one yields a model, which the interpreter must replay; an unknown answer or a
timeout proves nothing (invariant~I8). This section defines the queries for the core calculus
and proves that proving them all makes a node satisfy its contract.

\subsection{Facts, obligations and contexts}

Let $T$ be the standard model of the quantifier-free combination of bit-vectors, floating point
(with one NaN per sort), arrays and uninterpreted functions. The encoding of a node assigns a
term to every port, built bottom-up over symbols for the inputs, and records for each op $o_i$
a set of \emph{facts} $A_i$ (formulas the encoding may assume once $o_i$ has run) and a set of
\emph{goals} $G_i$ (formulas that must hold for $o_i$ to run). \Cref{tab:rules} lists them. A
fresh symbol is a new constant that occurs in no earlier formula; fresh symbols in a goal are implicitly
universally quantified.

\begin{table}[t]
\centering
\small
\renewcommand{\arraystretch}{1.25}
\begin{tabularx}{\linewidth}{@{}>{\raggedright\arraybackslash}p{3.1cm}>{\raggedright\arraybackslash}X>{\raggedright\arraybackslash}X@{}}
\toprule
Op $o_i$ & Goals $G_i$ & Facts $A_i$ \\
\midrule
primitive $f(\bar a)$ & $\mathit{dom}_f(\bar a)$ (absent when $f$ is total) &
  $p_i=\mathit{enc}_f(\bar a)$ \\
\op{CALL\_HASH}\newline$(g,\bar a)$ & $\Pre_g(\bar a)$ & $\Post_g(\bar a,p_i)$ \\
\op{ASSERT\_GUARD}\newline$(c,g_1,\bar a,g_2)$ & $c\Rightarrow\Pre_{g_1}(\bar a)$;\ \
  $\neg c\Rightarrow\Pre_{g_2}(\bar a)$ &
  $c\Rightarrow\Post_{g_1}(\bar a,r_1)$;\ \ $\neg c\Rightarrow\Post_{g_2}(\bar a,r_2)$;\ \
  $p_i=\mathrm{ite}(c,r_1,r_2)$ \\
\op{SIMD\_MAP}\newline$(g,xs,\bar b)$ & $0\le k<n\Rightarrow\Pre_g(xs[k],\bar b)$, with
  $n=\mathrm{len}(xs)$ and $k$ fresh & $\forall k<n.\ \Post_g(xs[k],\bar b,p_i[k])$, where $p_i$ has the
  shape of $xs$ \\
\op{GENERATE}\newline$(g,n,\bar b)$ & $0\le k<n\Rightarrow\Pre_g(k,\bar b)$, $k$ fresh &
  $\forall k<n.\ \Post_g(k,\bar b,p_i[k])$ \\
\op{FOLD\_BOUNDED}\newline$(g,xs,\bar b,\mathit{init})$,\newline invariant $I$ &
  $I(0,\mathit{init})$;\ \ $0\le k<n\wedge I(k,\mathit{acc})\Rightarrow
  \Pre_g(\mathit{acc},xs[k],\bar b)$;\ \ $0\le k<n\wedge I(k,\mathit{acc})\wedge
  \Post_g(\mathit{acc},xs[k],\bar b,r)\Rightarrow I(k+1,r)$; with $n=\mathrm{len}(xs)$ and
  $k,\mathit{acc},r$ fresh & $I(n,p_i)$ \\
effect op & none & environment facts (e.g.\ clock reads never decrease) \\
\bottomrule
\end{tabularx}
\caption{Goals and facts per op, following the specification~\cite{vgspec}. $p_i$ is the term
of the op's result: the encoding of a primitive op, and a fresh symbol for every other op.
$\Pre_g$ and $\Post_g$ are the callee's contract with its parameters and result replaced by the
given terms.}
\label{tab:rules}
\end{table}

A universally quantified fact is never given to a solver as such. The fresh-element rule of the
specification instantiates it at the index terms of the same collection that occur in the
query, in a fixed order, so the query stays quantifier-free. Goals with a \kw{forall} are
handled dually: the bound variable becomes a fresh constant.

\begin{definition}[Ordered contexts]\label{def:ordered}
Number the ops of a canonical node in canonical order. This is a topological order in which the
ports named by a fold's invariant also precede the fold, because canonical numbering treats
those ports as dependencies of the fold, as the prototype does. The context of op $o_i$ is
\[
\ctx_{<i} \;=\; \Pre\wedge\mathit{Box}\wedge\bigwedge_{j<i}A_j .
\]
The obligations of the node are $\ctx_{<i}\Rightarrow\gamma$ for every goal $\gamma\in G_i$,
and $\ctx_{<k}\Rightarrow\Post$ for the postcondition, where $k$ is the number of ops. An
encoding may drop conjuncts from a context, for instance those that share no symbol with the
goal.
\end{definition}

\begin{definition}[Verified node]\label{def:verified}
A node is \emph{verified} when its contract is non-vacuous ($\Pre\wedge\mathit{Box}$ is
satisfiable), every obligation is valid in $T$, as witnessed by an unsatisfiable query, and
every callee is verified, or is a registry binding whose evidence the caller's policy accepts.
\end{definition}

\begin{assumption}[Faithful encoding]\label{ass:faithful}
For every primitive op symbol $f$, the formula $\mathit{dom}_f$ defines $D_f$ in $T$, and the
term $\mathit{enc}_f$ denotes $\sem{f}$ on $D_f$. For contracts, formula terms denote their
meaning of \cref{sec:contracts}, with unspecified values as uninterpreted functions.
\end{assumption}

\Cref{ass:faithful} is the part of the verification-condition generator that must be trusted;
\cref{sec:float} discharges the delicate float cases, and \cref{thm:widths} below the integer
terms of formulas.

\subsection{Two logical lemmas}

\begin{lemma}[Generalization]\label{lem:gen}
Let $\Delta$ be a set of formulas in which the constant $k$ does not occur. If
$\Delta\wedge\rho(k)\wedge\neg\psi(k)$ is unsatisfiable, then
$\Delta\models\forall i.\ \rho(i)\Rightarrow\psi(i)$.
\end{lemma}
\begin{proof}
If a model $M$ of $\Delta$ had an element $i_0$ with $\rho(i_0)$ and $\neg\psi(i_0)$, then
reinterpreting $k$ as $i_0$ would leave $\Delta$ true, since $k$ does not occur in it, and
satisfy the query.
\end{proof}

\begin{lemma}[Weakening]\label{lem:weak}
If $\Delta'$ is a set of logical consequences of $\Delta$ and $\Delta'\wedge\neg\psi$ is
unsatisfiable, then $\Delta\models\psi$. In particular, instantiating a universally quantified
fact at any finite set of terms, or dropping conjuncts, is sound.
\end{lemma}
\begin{proof}
Every model of $\Delta\wedge\neg\psi$ is a model of $\Delta'\wedge\neg\psi$.
\end{proof}

The converse of \cref{lem:weak} fails, which is why a model of a weakened query may be
spurious: it satisfies the instances that were kept but may violate one that was dropped. The
specification reports such a model as \emph{unknown}, never as a counterexample
(\cref{thm:replay}).

\subsection{Soundness}

\begin{theorem}[Soundness of modular verification]\label{thm:soundness}
Let $N$ be a well-formed canonical node, let every callee of $N$ satisfy its contract, and
suppose every obligation of $N$ (\cref{def:ordered}) is valid in $T$. Under
\cref{ass:faithful}, $N$ satisfies its contract (\cref{def:satisfies}).
\end{theorem}
\begin{proof}
Fix an input $v$ with $\mathit{Box}(v)$, an interpretation $\iota$ with $\sem{\Pre}^\iota_v$,
and an admissible trace. We build, op by op, a $T$-model $M$ that interprets the input symbols
as $v$, the uninterpreted functions of unspecified values as $\iota$, and the symbols of each op
that has run as the values it computed. We prove by induction on $i\le k$:
\begin{center}
\emph{ops $o_0,\dots,o_{i-1}$ do not trap, and $M\models\ctx_{<i}$.}
\end{center}
For $i=0$, $\ctx_{<0}=\Pre\wedge\mathit{Box}$ holds by the choice of $v$ and $\iota$. For the
step, assume the claim for $i$. Each goal $\gamma\in G_i$ satisfies $\ctx_{<i}\Rightarrow\gamma$
in $T$, with its fresh symbols universally quantified (\cref{lem:gen}); interpreting those
symbols as any values keeps $M\models\ctx_{<i}$, since they do not occur there, so $\gamma$
holds in $M$ for every such choice. We show that $o_i$ does not trap and that $A_i$ can be made
true by interpreting the fresh symbols of $A_i$, which occur in no earlier formula.
\begin{itemize}
\item \emph{Primitive op.} If $f$ is partial, its goal $\mathit{dom}_f(\bar a)$ holds of the
  operand values, so $\bar a\in D_f$ by \cref{ass:faithful}, and $o_i$ does not trap. Setting
  $p_i$ to $\sem{f}(\bar a)$ satisfies $p_i=\mathit{enc}_f(\bar a)$.
\item \emph{Call.} $\Pre_g(\bar a)$ holds, so the callee, which satisfies its contract, does not
  trap and returns a result satisfying $\Post_g$; interpret the fresh $p_i$ as that result.
\item \emph{Guard.} Exactly one branch runs. If $c$ holds, $\Pre_{g_1}(\bar a)$ holds, $g_1$
  returns $r_1$ with $\Post_{g_1}$, and we interpret $r_2$ arbitrarily, which satisfies the
  implication guarded by $\neg c$; symmetrically if $c$ fails. Then $p_i=\mathrm{ite}(c,r_1,r_2)$
  is the value computed.
\item \emph{Map and generate.} Instantiating the fresh $k$ at each index shows
  $\Pre_g(xs[k],\bar b)$ for every element, so no kernel application traps, and each result
  satisfies $\Post_g$; interpret $p_i$ as the computed collection. Every finite set of
  instances of the universal fact then holds (\cref{lem:weak}).
\item \emph{Fold.} Let $\mathit{acc}_0=\mathit{init}$ and
  $\mathit{acc}_{q+1}=g(\mathit{acc}_q,xs[q],\bar b)$. By induction on $q\le n$, $I(q,
  \mathit{acc}_q)$ holds and steps $0,\dots,q-1$ do not trap: the case $q=0$ is the first goal;
  for $q<n$, instantiating $k,\mathit{acc}$ as $q,\mathit{acc}_q$ in the second goal gives
  $\Pre_g(\mathit{acc}_q,xs[q],\bar b)$, so the kernel does not trap and its result
  $\mathit{acc}_{q+1}$ satisfies $\Post_g$; instantiating $r$ as $\mathit{acc}_{q+1}$ in the
  third goal gives $I(q+1,\mathit{acc}_{q+1})$. Interpret $p_i$ as $\mathit{acc}_n$, so
  $I(n,p_i)$ holds. The invariant therefore holds at every step.
\item \emph{Effect op.} It does not trap, and an admissible trace satisfies the environment
  facts.
\end{itemize}
Hence $M\models\ctx_{<i+1}$. For $i=k$, no op traps, so by \cref{lem:unr-sem} the outputs of $N$
are the values of the output terms in $M$, and the postcondition obligation gives
$\sem{\Post}^\iota$ of the result.
\end{proof}

Applied bottom-up along the acyclic call graph, \cref{thm:soundness} shows that every node of a
verified program, all of whose registry bindings have \emph{proved} evidence, satisfies its
contract. The guarantee is relative to the trusted computing base: the canonicalizer, the type
checker, the verification-condition generator (\cref{ass:faithful}), the solvers, and, for
compiled code, the code generator, LLVM, the runtime and the operating system.

\begin{corollary}[No local vacuity]\label{cor:local-vacuity}
Under the hypotheses of \cref{thm:soundness}, if $\Pre\wedge\mathit{Box}$ is satisfiable, then
so is every context $\ctx_{<i}$.
\end{corollary}
\begin{proof}
For any input satisfying $\Pre\wedge\mathit{Box}$, the model built in the proof satisfies
$\ctx_{<i}$.
\end{proof}

So once the contract itself is shown non-vacuous, no obligation can be proven merely because
the facts of earlier ops contradict each other; this would fail for registry bindings whose
evidence is only \emph{tested}, whose postconditions may be false.

\subsection{The ordering proviso is necessary}

\Cref{def:ordered} lets an op's goals see only the facts of earlier ops. The specification
describes contexts as ``the caller's facts'' and keeps ``only assumptions that share symbols
with the goal, closed under sharing''; it does not exclude an op's own facts or those of later
ops. That omission is unsound.

\begin{lstlisting}[float=t,caption={A program that traps at $x=5$ but whose obligations all
become valid if contexts may contain an op's own facts.},label=lst:ordering]
NODE @g {
  SIGNATURE: (y: i32) -> (i32);
  EFFECTS: [];
  CONTRACT {
    PRE:  y != 0;
    POST: y != 0 && result == 1;
  }
  BODY {
    RETURN(1:i32);
  }
}

NODE @f {
  SIGNATURE: (x: i32) -> (i32);
  EFFECTS: [];
  CONTRACT {
    POST: true;
  }
  BODY {
    %0 = SUB.wrap(x, 5:i32);
    %1 = DIV.checked(100:i32, %0);
    %2 = CALL_HASH(@g, %0);
    %3 = ADD.wrap(%1, %2);
    RETURN(%3);
  }
}
\end{lstlisting}

\begin{proposition}[Unordered contexts are unsound]\label{prop:ordering-needed}
In \cref{lst:ordering}, \kw{@g} satisfies its contract and is non-vacuous, and \kw{@f} traps at
$x=5$. If each obligation of \kw{@f} may use the facts of every op that shares a symbol with
its goal, then all obligations of \kw{@f} are valid.
\end{proposition}
\begin{proof}
\kw{@g} returns $1$, so its postcondition holds whenever its precondition does, and $y=1$
satisfies the precondition. Write $p_j$ for the port \kw{\%j} of the listing. In \kw{@f},
$p_0=x-5$ in wrapping \kw{i32} arithmetic; at $x=5$,
$p_0=0$ and \op{DIV.checked} traps (so does the call, whose precondition fails). The goals of
\kw{@f} are $p_0\neq0$ for \op{DIV.nonzero}, a trivially true \op{DIV.no\_overflow} (the
dividend is $100$), and $p_0\neq0$ for \op{CALL.pre.0}. The fact of the call is
$\Post_g(p_0,p_2)=(p_0\neq0\wedge p_2=1)$, which shares the symbol $p_0$ with both goals. With it in
the context, both goals follow immediately; the postcondition \kw{true} is valid. So every
obligation is valid although \kw{@f} traps.
\end{proof}

Under \cref{def:ordered} the call's own fact is excluded from the context of its goal, the goal
$p_0\neq0$ fails, and \kw{@f} is rejected, as it must be. The listing parses and assembles with
the prototype's toolchain (\cref{app:repro}). Guarding call facts by the callee's precondition,
$\Pre_g(\bar a)\Rightarrow\Post_g(\bar a,p_i)$, would also reject this program, but it does not
suffice for folds: an invariant may name a port of the body, and a fold whose own fact
$I(n,p_i)$ entered the context of its own goals could prove its base case from its conclusion.
Ordering contexts by a topological order handles every op uniformly, and canonical numbering
already provides one.

\subsection{Exact integers in bit-vectors}

Formula integers are mathematical integers, but every query is quantifier-free bit-vector
arithmetic, which all three solvers accept and which is decidable, whereas integer arithmetic
with products of variables is not~\cite{matiyasevich1970}. The specification encodes each
integer term $t$ as a signed bit-vector of a width $w(t)$ chosen so that nothing wraps.

\begin{definition}[Widths]\label{def:widths}
$w$ is defined on integer terms by: $w(x)=N+1$ for an input of type $\mathsf{u}N$ or
$\mathsf{i}N$; $w(\mathrm{len}(xs))=65$; $w(c)=\mathrm{bits}(|c|)+1$ for a literal $c$, where
$\mathrm{bits}(m)$ is the number of binary digits of $m$ ($\mathrm{bits}(0)=0$); for a variable
bound by $\kw{forall}(i\ \kw{in}\ a..b)$, $w(i)=\max(w(a),w(b))$;
$w(a\pm b)=\max(w(a),w(b))+1$; $w(a\cdot b)=w(a)+w(b)$; $w(-a)=w(\mathrm{abs}(a))=w(a)$;
$w(a/b)=w(a)$; $w(a\mathbin{\%}b)=\min(w(a),w(b))$; $w(\min(a,b))=w(\max(a,b))=\max(w(a),w(b))$.
An operator sign-extends its operands to the largest of their widths and its result width,
computes there with \kw{bvadd}, \kw{bvsub}, \kw{bvmul}, \kw{bvneg}, \kw{bvsdiv} or
\kw{bvsrem}, and truncates to its result width; a comparison sign-extends both sides to the
wider width. Division and remainder by zero are uninterpreted functions of their operands.
\end{definition}

\begin{theorem}[Exact integer encoding]\label{thm:widths}
For every integer term $t$ and every assignment of in-range values to its free variables under
which every range variable satisfies its bounds and every divisor is non-zero, the bit-vector
value of $t$, read as a signed integer of width $w(t)$, equals the mathematical value
$\sem{t}$, and $|\sem{t}|<2^{w(t)-1}$. Every comparison of integer terms is therefore decided
exactly.
\end{theorem}
\begin{proof}
By induction on $t$. Write $W$ for the computation width of an operator and recall that a
$W$-bit signed bit-vector holds exactly the integers in $[-2^{W-1},2^{W-1}-1]$, so a value of
magnitude below $2^{W-1}$ is held without wrapping and is never the most negative value.
\begin{itemize}
\item Inputs: a $\mathsf{u}N$ value lies in $[0,2^N-1]$ and an $\mathsf{i}N$ value in
  $[-2^{N-1},2^{N-1}-1]$; both have magnitude below $2^N=2^{w-1}$, and zero- or sign-extension
  by one bit preserves the value. $\mathrm{len}$ is below $2^{64}=2^{65-1}$. A literal satisfies
  $|c|<2^{\mathrm{bits}(|c|)}=2^{w(c)-1}$.
\item A range variable satisfies $a\le i<b$, so $-2^{w(a)-1}<i<2^{w(b)-1}$, and
  $|i|<2^{w(i)-1}$. Where the bounds fail, $i$ occurs only under the guard
  $a\le i<b$, which is evaluated exactly, so its value there does not matter.
\item $a\pm b$: $W=w(a\pm b)$ and $|a\pm b|\le|a|+|b|<2^{w(a)-1}+2^{w(b)-1}\le2^{W-1}$.
\item $a\cdot b$: $W=w(a)+w(b)$ and $|ab|<2^{w(a)-1}2^{w(b)-1}=2^{W-2}$.
\item $-a$ and $\mathrm{abs}(a)$: the magnitude is $|a|<2^{w(a)-1}$, so \kw{bvneg} is exact.
\item $a/b$ with $b\neq0$: $W=\max(w(a),w(b))$; since $|a|<2^{w(a)-1}\le2^{W-1}$, $a$ is not
  the most negative $W$-bit value, so \kw{bvsdiv}, which truncates toward zero as \op{DIV} does,
  is exact, and $|a/b|\le|a|<2^{w(a)-1}$ survives truncation to $w(a)$ bits.
\item $a\mathbin{\%}b$ with $b\neq0$: \kw{bvsrem} takes the sign of the dividend, as \op{MOD}
  does, and is exact in width $W$; $|a\mathbin{\%}b|\le|a|<2^{w(a)-1}$ and
  $|a\mathbin{\%}b|\le|b|-1<2^{w(b)-1}$, so the value fits $\min(w(a),w(b))$ bits.
\item $\min$ and $\max$ return an operand, of magnitude below $2^{\max(w(a),w(b))-1}$.
\end{itemize}
A comparison extends both sides to a common width that holds both values, where the signed
comparison is the integer comparison.
\end{proof}

A width that is too narrow would be a soundness bug: a wrapped sum can make a false
postcondition look proven. The independent reviews of this paper's proofs found no term on which
the widths fail (\cref{sec:status}).

\subsection{Folds and specifications inside formulas}

A postcondition such as \kw{result == spec:sha256(arg0)} needs a way to say what a fold has
computed. The specification adds two terms. $\kw{fold:}f(\mathit{init},xs,k,\bar b)$ is the
left fold of kernel $f$ over the first $k$ elements; the solver sees an uninterpreted function
$F$ and receives, for each ground application, the two instances
\begin{gather*}
k=0\Rightarrow F(\mathit{init},xs,k,\bar b)=\mathit{init},\\
0<k\le\mathrm{len}(xs)\Rightarrow F(\mathit{init},xs,k,\bar b)=
f(F(\mathit{init},xs,k-1,\bar b),xs[k-1],\bar b),
\end{gather*}
never chained further. $\kw{spec:}s(\bar t)$ is an uninterpreted function $S$ for the executable
reference node $s$, unfolded one level.

\begin{proposition}[Fold and spec instances are sound]\label{prop:foldspec}
Adding any set of the fold instances above is sound. Adding, for an application of
$\kw{spec:}s$, the guarded facts
$\Pre_s(\bar t)\Rightarrow S(\bar t)=\mathit{body}_s(\bar t)$ and
$\Pre_s(\bar t)\Rightarrow\Post_s(\bar t,S(\bar t))$ is sound when $s$ satisfies its contract.
Adding $\Pre_s(\bar t)$ unguarded is unsound, and so is adding $\Post_s(\bar t,S(\bar t))$
unguarded, since a postcondition may restate the precondition.
\end{proposition}
\begin{proof}
The intended interpretation of $F$, the true fold for $0\le k\le\mathrm{len}(xs)$ and anything
for other $k$,
satisfies both instance schemes for all arguments, and it is one of the interpretations over
which validity quantifies; so every instance is a consequence of the intended meaning, and
\cref{lem:weak} applies. The same holds for $S$ interpreted as the function computed by $s$ on
its precondition: on $\Pre_s$, $s$ does not trap, returns $\mathit{body}_s(\bar t)$ and
satisfies $\Post_s$. For the last claim, let $s$ return an integer and have
$\Pre_s(y)\equiv y>0$, and give a node with precondition \kw{true} the postcondition
$\Post\equiv x>0\wedge(x>0\Rightarrow\kw{spec:}s(x)=\kw{spec:}s(x))$, in which the
\kw{spec:} term is well defined because it is guarded by $x>0$. The postcondition is false at
$x=0$, yet the unguarded fact $\Pre_s(x)$, that is $x>0$, makes its obligation valid. The same
happens with an unguarded $\Post_s$ if $s$ has the postcondition $y>0$, which $s$ satisfies.
\end{proof}

The specification says that the precondition and postcondition of a \kw{spec:} call ``are added
as facts''; \cref{prop:foldspec} shows they must be added under the guard $\Pre_s(\bar t)$, and
the well-definedness obligations of \cref{sec:contracts} then establish that guard wherever
the term is read. The independent soundness review reached the same conclusion
(\cref{sec:status}).

The natural invariant of a fold, $\mathit{acc}=\kw{fold:}f(\mathit{init},xs,k,\bar b)$,
becomes provable with one instance per query: the base case uses the instance at $k=0$, and the
step goal $r=F(\dots,k+1,\dots)$ follows from the instance at $k+1$ and the hypothesis
$\mathit{acc}=F(\dots,k,\dots)$, provided the step query relates $r$ to
$f(\mathit{acc},xs[k],\bar b)$ exactly. That holds when the kernel is inlined in the step query,
as it is inside formulas; if the step used only $\Post_f$, the natural invariant would be
provable exactly when $\Post_f$ determines the result. Either way soundness is unaffected.

\subsection{Non-vacuity}

\begin{lemma}[Ex falso]\label{lem:exfalso}
If $\Pre\wedge\mathit{Box}$ is unsatisfiable, every obligation of the node is valid.
\end{lemma}
\begin{proof}
Every context contains $\Pre\wedge\mathit{Box}$.
\end{proof}

A node with an unsatisfiable precondition would therefore verify whatever its body does, which
is why \cref{def:verified} requires non-vacuity. The specification establishes it by an input
from the concrete pre-screen, or by a solver model that the interpreter replays; either is a
certificate of satisfiability that can be checked by evaluation, and a check that ends in
\emph{unknown} rejects the node.

\subsection{Counterexamples and their replay}

Modular verification abstracts on purpose: a callee is replaced by its contract, a \kw{spec:}
or \kw{fold:} application by an uninterpreted function, and a universal fact by some of its
instances. A model of a query may therefore choose values that no real execution produces. The
specification replays every model twice: \emph{model-guided} replay runs the interpreter with
the model's inputs and with the model's values for every abstracted term, and checks every
assumption in full; \emph{concrete} replay runs the real callees.

\begin{assumption}[Interpreter fidelity]\label{ass:interp}
For any inputs and any values assigned to the abstracted terms, the interpreter computes, for
every port, the value that the port's term denotes under that assignment.
\end{assumption}

\begin{theorem}[Replay classification]\label{thm:replay}
Let $M$ be a model of an obligation query $\Delta\wedge\neg\gamma$, and suppose every
assumption of the node, including every universal fact in full, holds in $M$.
\begin{enumerate}[label=(\alph*)]
\item Under \cref{ass:faithful,ass:interp}, model-guided replay reproduces the violation of
  $\gamma$. Hence a model-guided replay that does not reproduce it reveals an unfaithful encoding,
  a verifier bug.
\item If concrete replay reproduces the violation, the model's inputs satisfy the precondition
  and the box, and the node, run on them, violates $\gamma$: a genuine counterexample.
\item If model-guided replay reproduces the violation but concrete replay does not, then
  $\gamma$ does not follow from the precondition and the callees' contracts: a contract gap.
\end{enumerate}
If some universal fact that was only partially instantiated fails in $M$, $M$ is evidence of
nothing, and the verdict is unknown.
\end{theorem}
\begin{proof}
(a) Model-guided replay fixes the inputs and the abstracted terms to their values in $M$; by
\cref{ass:faithful,ass:interp} every port then takes its value in $M$, so $\gamma$ evaluates as
in $M$, to false. (b) Concrete replay evaluates the precondition and $\gamma$ on the real
execution; a reproduction is a real input satisfying the precondition on which $\gamma$ fails.
(c) $M$ satisfies every assumption, among them all that the encoding knows about the callees,
namely their contracts, and falsifies $\gamma$; so $\gamma$ is not a consequence of them. For
the last claim, a model that violates an assumption is not a model of the full set of
assumptions, so it neither refutes $\gamma$ nor reveals a bug.
\end{proof}

Only an unsatisfiable query proves an obligation. Counterexample, contract gap, unknown and
error are all failures, distinguished only to tell the agent which repair applies: fix the
body; strengthen a callee's postcondition or the node's precondition; split the obligation or
guard the operation; or, for an error, stop, since the verifier itself is wrong.

\subsection{Costs}

The specification proves \kw{COST} bounds in an abstract model: one unit per executed op; a
map, fold or generate adds $n$ times its kernel's cost; \op{ASSERT\_GUARD} adds the larger of
its branches; \op{CALL\_HASH} adds the callee's proven bound.

\begin{proposition}[Cost bounds]\label{prop:cost}
For every input on which a node does not trap, the number of op executions, counted through
kernels and callees, is at most the computed cost expression. Equality holds when the node and
its callees contain no \op{ASSERT\_GUARD} and every callee's bound is exact.
\end{proposition}
\begin{proof}
By induction on the call graph and the op order: straight-line ops contribute exactly one each;
a map over $n$ elements executes its kernel $n$ times, each costing its kernel's constant cost
(kernels are fixed-size); a guard executes one branch, whose cost is at most the larger; a call
executes at most its callee's proven bound.
\end{proof}

The specification calls the cost expression exact; \cref{prop:cost} shows it is an upper bound
in general, exact only without guards. Arena usage is bounded the same way.

\subsection{The trust boundary}

Compiled code assumes preconditions: proven facts become \kw{nsw}, \kw{inbounds} and
\kw{llvm.assume}. A caller written in C or Rust is not verified, so every exported node gets a
checked entry point that evaluates the precondition on the actual inputs, rejects Boolean
bytes other than $0$ and $1$ and non-canonical NaN patterns, checks the arena size, sets the
default floating-point environment, and only then calls the internal function. When the checked
entry returns success, the hypotheses of \cref{thm:soundness} hold for that call, so the proof
facts are never applied to inputs that falsify them.
